% !PRINTED 111
\weberpage{111}
%<<CONT
und ferner die Summe

\weq{4}{x_0 + x_1 + x_2 + x_3 = 4A}

rationale Zahlen.

Wenn wir aus (2) die Quadratwurzeln ausziehen und berücksichtigen, dass sich
$\sqrt{a}$ nach (3) von $\sqrt{bc}$ nur durch einen rationalen Factor
unterscheidet, so folgt durch Addition:

\weq{5}{x_0 = A + B\sqrt{b} + C\sqrt{c} + D\sqrt{bc},}

worin $A, B, C, D$ rationale Zahlen sind; und daraus erhält man $x_1, x_2,
x_3$, wenn man die Vorzeichen von $\sqrt{b}, \sqrt{c}$ ändert.

Nach Band I, \S.\,171 können aber alle Quadratwurzeln, nöthigenfalls unter
Zuziehung von $i$, was ja selbst eine Kreistheilungszahl ist, rational durch
die Kreistheilungsperioden (die \gesperrt{Gauss}'schen Summen) ausgedrückt
werden, so dass also in (5) schon der Beweis unseres Satzes liegt.

\noindent\hrulefill

Ist die Gruppe der biquadratischen Gleichung cyklisch, so können wir die
Wurzeln so anordnen, dass die Gruppe aus den Permutationen

\weq{6}{1, \ \ (0,1,2,3), \ \ (0,2)(1,3), \ \ (0,3,2,1)}

besteht, und dann ist zu setzen:

\weq{7}{\begin{aligned}
4A &= x_0+x_1+x_2+x_3 \\
(i,x_0) &= x_0+ix_1-x_2-ix_3 \\
(-1,x_0) &= x_0-x_1+x_2-x_3 \\
(-i,x_0) &= x_0-ix_1-x_2+ix_3;
\end{aligned}}

darin ist $A$ und $(-1,x_0)^2=m$ rational, und daher kann $(-1,x_0)$ durch
Kreistheilungszahlen ausgedrückt werden. Die vierten Potenzen

\weq{8}{(i,x_0)^4 = a+bi, \quad (-i,x_0)^4 = a-bi}

sind Zahlen des Körpers $R(i)$, und um dieselben Schlüsse wie bei den
cubischen Gleichungen ziehen zu können, kommt es also nur noch darauf an,
$a+bi$ und $a-bi$ als vierte Potenzen von Kreistheilungszahlen darzustellen.
Dazu dient noch der Satz, dass

\weq{9}{(i,x_0)(-i,x_0) = c}

eine rationale Zahl ist.
